\section{这本书该怎么读}

\label{sec:0-3}

这本书当然可以从头顺读，但未必需要从头顺读。因为它同时服务三种常见读者：想把理论基础补扎实的纯数学读者，已经懂 Boyd 母语、主要关心算法接口的读者，以及更在意建模闭环和应用落点的管理科学读者。不同入口并不意味着不同版本的书，只是意味着你先回答哪个问题。

如果你主要关心纯数学基础，那么最自然的路径是：第 \ref{chap:01} 章建立拓扑与描述集合论语言，第 \ref{chap:02} 章进入测度与随机空间，第 \ref{chap:03} 章和第 \ref{chap:04} 章建立泛函分析与分离工具，再进入第 \ref{chap:05} 章和第 \ref{chap:06} 章的次微分与对偶理论。这条线最慢，但也最稳，因为它把后文一切“算法为何有定义”“乘子为何存在”“弱极限为何足够”这些问题都提前补齐。

如果你主要关心算法，那么不必在一开始就吞下全书。更短的一条路是：先把第 \ref{chap:03} 章中的 Hilbert 空间、弱拓扑和谱论读到能用，再接第 \ref{chap:05} 章的次微分、第 \ref{chap:09} 章的单调算子与预解式，最后进入第 \ref{chap:10} 章的前向--后向分裂、Douglas--Rachford 和 ADMM。沿这条路线读时，Boyd 会一直作为有限维母语出现，但你会逐步看清：这些算法真正处理的不是矩阵，而是算子零点与原--对偶系统。

如果你主要关心管理科学、控制或应用建模，那么更合适的入口是先读随机过程与对偶接口，再回补基础。具体说，可以从第 \ref{sec:2-4} 节的随机分析入口读起，接着读第 \ref{sec:6-4} 节 Fenchel--Rockafellar 对偶、第 \ref{chap:07} 章锥优化与 KKT，再看第 \ref{chap:11} 章和第 \ref{chap:12} 章的控制、交易、图网络与推荐系统。这样你会更快看到闭环：为什么约束资格、风险对偶、分布式一致性和表示匹配都能被放回同一套凸优化语言。

无论从哪条路切入，建议都一样：把 Boyd 留在心里，但不要把它当成天花板。有限维模型是最重要的直觉来源，却不是全部世界。你越早接受“变量可以是函数、测度或轨道”，后文的拓扑、对偶和算子就越不再像门槛，而像一条自然延长线。
